\section[PHT3D Exercise 4: Transport and nitrification of ion exchangeable ammonium]{PHT3D Exercise 4: Transport and nitrification of ion exchangeable ammonium}

This modelling example is based on a field site contamination problem near
Mansfield UK \citep[see, e.g.,][]{broholm98, phddavison,jones1998, davison2000,
jones2001}, where ammonium liquor, a by-product of the production
of smokeless fuel, has polluted groundwater over several decades.
A reactive transport modelling study that integrated 
and reproduced the major processes that were believed to occur at the
field site was presented by \cite{haerens_phd}. 
One of the key features observed at the site is the
strongly retarded migration of ammonium and the geochemical footprint that
was left behind as a result of the cation exchange of ammonium.  
In the modelling example the development of an ammonium plume
and the subsequent flushing of ammonium contaminated groundwater by
pristine background water is simulated.  The processes included in
the example are advection, dispersion, cation exchange
and the kinetically controlled oxidation of ammonium.
Dispersion, ion exchange and nitrification act as attenuation 
processes for ammonium.

For simplicity the two-dimensional reactive transport problem is set up
for the same model domain as the previous exercise.
However, the simulation period is now divided into two different phases. 
The first phase represents the period 
of active contamination during which the plume grows successively
while the second phase represents the period after the source 
was exhausted. During this second phase clean groundwater will enter 
the model domain over the entire depth of the upstream boundary. 
Furthermore, we also include now the simulation of groundwater recharge
in this model.

\subsection[Modification of the flow problem]{Modification of the flow problem}

For convenience the flow model is similar to the one used in the previous exercise, i.e,
the spatial dimensions were modified from the original field setting.
The details of the flow model and its discretisation are listed in \ref{bruno_flow}.
To adapt the MODFLOW model from the previous exercise to the present case, 
proceed as follows:

\begin{itemize}

\item Make a copy of the flow model that was readily prepared for this exercise
and save it in a new folder.
The model files are located in the 
folder {\color{red} \textbf{2d\_flow\_problem\_haerens}}.

\item Select {\color{blue}Parameters} and select the {\color{blue}Time} button.
and change the {\color{blue}Total Simulation Time} to 3000 days and select 
a {\color{blue}Step Size} to 40 days. 

\item To add groundwater recharge we first need to activate the 
MODFLOW recharge package such that it is availale within ORTI.
To do this, go to ORTI's top menu and select {\color{blue}Add-in} 
$\rightarrow$ {\color{blue}Modflow Modules}. You will see a list of 
all available MODFLOW packages and you will see which ones are already activated.
Add the recharge package by ticking the box near {\color{blue}RCH}.
Now go to {\color{blue}Spatial Attributes} $\rightarrow$ {\color{blue}MODFLOW} $\rightarrow$ {\color{blue}RCH} $\rightarrow$ {\color{blue}rch.2 Recharge Value} to specify a rechare rate 
of 0.001 $m/day$, i.e., 1 mm per day. Don't forget to press {\color{blue}'OK'} 
once you have entered the value.   

\item Now rerun MODFLOW to regenerate the flow-file (mt3d.flo), which will later be used 
by the transport model.

\end{itemize}

\begin{table}
\caption{Flow and transport parameters used for PHT3D Exercise 4.}
\begin{center}
\begin{tabular}{lc}
\toprule
Flow simulation                     & steady state \\
Total simulation time $(days)$       & 3000 \\
Time step size                       & 40 \\
Model \  length\ $ (m)      $         &  100            \\
Model \  thickness\ $ (m)      $         &  10            \\
Grid \ spacing\ $\Delta x $ (m)                &    4            \\
Grid \ spacing\ $\Delta z $ (m)                &    0.5            \\
Porosity\     $            $         &  0.35           \\
Horizontal hydraulic conductivity\ $(m)           $         &  1 m/d         \\ 
Vertical   hydraulic conductivity\ $(m)           $         &  1 m/d         \\ 
Piezometric head upstream boundary\ $(m)          $         &  12 m          \\ 
Piezometric head downstream boundary\ $(m)        $         &  10 m          \\ 
Longitudinal dispersivity\ $(m)         $         &  2.5          \\ 
Transverse vertical dispersivity\ $(m)         $         &  0.025          \\ 
\bottomrule

\end{tabular}
\label{bruno_flow}
\end{center}
\end{table}


\subsection[Data input for solute transport properties]{Data input for solute transport properties}

For the simulation of the solute transport of the mobile components the model parameters that control advective and dispersive transport need to be entered. 

\begin{itemize}

\item Go to {\color{blue}Parameters} $\rightarrow$ {\color{blue}Transport} and select the {\color{blue}Parameter} button

\item In the dialog box that subsequently appears, select {\color{blue}ADV} and select under {\color{blue}adv.1 Solver Parameters} the {\color{blue}3rd-order TVD Scheme (ULTIMATE)} as the Solution Scheme. Leave the default parameters and click {\color{blue}OK}.

\item Select {\color{blue}Spatial Attributes} $\rightarrow$ {\color{blue}MT3DMS} $\rightarrow$ {\color{blue}DSP}, enter a longitudinal dispersivity value of 2.5 m. and click {\color{blue}OK}

\item Select {\color{blue}Spatial Attributes} $\rightarrow$ {\color{blue}MT3DMS} $\rightarrow$ {\color{blue}DSP}, and specify that the ratio between transverse vertical dispersivity and the longitudinal dispersivity is 0.01.
This will define that the actual value is 0.025 m (= 25 mm). 

\item Make a quick test and assess the conservative transport behaviour in this setting.
To do this select {\color{blue}Spatial Attributes} $\rightarrow$ {\color{blue}MT3DMS} $\rightarrow$ {\color{blue}BTN} $\rightarrow$ {\color{blue}btn.13 Concentrations} and specify 
a small zone in the upper left corner of the model (near the inflow boundary) 
for which a concentration of 1 (or any other concentration > 0) is defined.

\item Under {\color{blue}Parameters} $\rightarrow$ {\color{blue}Transport} use the {\color{blue}Write Button} to write all the input files required to run MT3DMS

\item Under {\color{blue}Parameters} $\rightarrow$ {\color{blue}Transport} use the 
{\color{blue}Plume} button to start the MT3DMS simulation

\item Inspect the results by navigating in the {\color{blue} Results} panel of ORTI.
For example to see the concentration contours select {\color{blue}Results} $\rightarrow$ 
and select {\color{blue}Tracer}. Then select the first timestep 
under $\rightarrow$ {\color{blue}Tstep} and observe the plume spreading by using the 
right arrow button to move to subsquent timesteps.  

\end{itemize}

\subsection[Reactive transport]{Reactive transport}

Based on the flow and conservative transport model the next step is to define 
the input for the reactive transport.

\begin{itemize}

\item Copy the database file that will used for this PHT3D exercise
into the folder that contains all files for this exercise and make sure
its name is {\color{red} \textbf{pht3d\_datab.dat}}.
You can use the same database as in PHT3D Exercise 1, i.e., the standard PHREEQC database.

\item Go to {\color{blue}Parameters} $\rightarrow$ {\color{blue}Chemistry} and select the
{\color{blue}Import database} button. ORTI will then read and interpret 
the PHT3D database file.

\item In this next step we define all initial concentrations of all aqueous components and 
also the two types of water compositions at the upstream model boundary 
(which will be allocated to the inflowing groundwater). 
Activate the equilibrium components (tick the checkboxes) that are included
and also enter the concentrations provided in \ref{bruno_caqu}. 
Note, that there is no need to include aqueous components that will not play a role
in the simulations, such as {\color{red} \textbf{Mn(2)}}, {\color{red} \textbf{Mn(3)}} and  {\color{red} \textbf{Si}}. However, consider and activate all valence
states of redox-sensitive components. For example, activate all valence states of nitrogen,
i.e., {\color{red} \textbf{N(5)}}, {\color{red} \textbf{N(3)}}, {\color{red} \textbf{N(0)}} and {\color{red} \textbf{Amm}}. 

\item The composition of the ambient water (initial water composition at the start of the simulation) 
needs to be entered under the {\color{blue}Solutions} Tab 
in the {\color{blue}Background} column.
This solution can also be used to define the recharge water composition.
The water composition of the contaminated water will be entered in the 
Tab for {\color{blue}Solu 1}. If the concentration of
an aqueous component is 0 it is not necessary to enter the value.

\item Calcite is the only mineral that is considered in this simulation.
It must be activated and the corresponding initial concentrations must be entered 
under the {\color{blue}Phase} Tab.

\item The allocation of the appropriate water compoitions as model boundary condition can be
done by selecting {\color{blue}Spatial Attributes} $\rightarrow$ {\color{blue}MT3DMS} $\rightarrow$ {\color{blue}BTN} $\rightarrow$ {\color{blue}btn.12 Boundary Condition} and defining a zone (line)
for the lower section of the inflow (upstream) boundary (between y = 0m and y = 7m)
for which a value of -1 is allocated. This defines that the concentrations at these grid cells 
will not change during the entire simulation. The concentrations that ORTI will allocate to the 
model's grid cells are the ones defined as initial concentrations. 
This means that the water composition that was defined as background
water composition will at these locations also be used as inflow water composition. 

\item To additionally define that the two different water compositions
that sequentially enter the boundary's upper section we need to 
select {\color{blue}Spatial Attributes} $\rightarrow$ {\color{blue}PHT3D} $\rightarrow$ {\color{blue}PH} $\rightarrow$ {\color{blue}ph.4 Source \/ Sink} and define a zone (line) for the uppermost 3 m (y= 7m to y = 10m)
of the boundary. In the dialog that comes up with the {\color{blue}Zone} selection
enter two lines under {\color{blue}Zone Value}.
In the first line enter $0\ 1$ and in a separate second line enter
$1000\ 0$. This means that $Solution 1$ (as defined in the {\color{blue}Solutions} Tab) 
will be used until day 1000 and $Solution\ 0$ (which is also the background water composition) 
will be used until day 3000.

\item To correctly define the recharge water composition 
select {\color{blue}Spatial Attributes} $\rightarrow$ {\color{blue}PHT3D} $\rightarrow$ {\color{blue}PH} $\rightarrow$ {\color{blue}ph.5 Recharge} and define a zone (line) across the uppermost layer. 
In the dialog that comes up with the {\color{blue}Zone} selection
enter 0 under {\color{blue}Zone Value} to define that $Background\ Solution$ 
will also be used to define the recharge water composition.

\end{itemize}

The pollution source will be active during the first 1000 days, while 
during the 2000 days following the pollution event the inflowing water is similar 
to the composition of the ambient water.


\begin{table}
\caption{Concentrations of aqueous components and initial mineral concentration in PHT3D Exercise 4.}
\begin{center}
\begin{tabular}{l c c c c }
\toprule
      & \multicolumn{1}{c}{Background and} & \multicolumn{1}{c}{Contaminated water} \\
      & \multicolumn{1}{c}{flushing water} &                & \\
      & \multicolumn{1}{c}{$C_{backgr}$, $C_{flush}$} & \multicolumn{1}{c}{$C_{cont.}$} \\
      &    $mol\ l^{-1}$   & $mol\ l^{-1}$ \\ 
\midrule
O(0)     & 2.51 $\times$ $10^{-4}$    & $0$  \\
Na       & 8.62 $\times$ $10^{-4}$    & 1.30 $\times$ $10^{-3}$ \\
K        & 1.24 $\times$ $10^{-4}$    & 1.30 $\times$ $10^{-4}$ \\
Ca       & 1.83 $\times$ $10^{-3}$    & 1.50 $\times$ $10^{-4}$ \\
Mg       & 1.38 $\times$ $10^{-3}$    & 5.00 $\times$ $10^{-5}$ \\
Amm      & $0$                        & 6.87 $\times$ $10^{-3}$ \\
Cl       & 1.74 $\times$ $10^{-3}$    & 3.23 $\times$ $10^{-3}$ \\
S(6)     & 9.89 $\times$ $10^{-4}$    & 1.56 $\times$ $10^{-3}$ \\
N(5)     & 8.88 $\times$ $10^{-4}$    & $0$ \\
N(3)     &                            & $0$ \\ 
N(0)     &    $0$                     & $0$ \\ 
C(4)     & 2.82 $\times$ $10^{-3}$    & 2.92 $\times$ $10^{-3}$ \\ 
C(-4)    &  $0$                       &   $0$ \\ 
 & & \\
pH    & 7.9               & 8.3             \\
pe    & 13.5              & 0            \\
\midrule 
        &      $C_{init}$       &          \\
                &        $mol\ l_b^{-1}$     &          \\
Calcite &      0.1             &          \\
\bottomrule

\end{tabular}
\label{bruno_caqu}
\end{center}
\end{table}


\subsection[Running PHT3D]{Running PHT3D}

To run PHT3D, proceed as earlier in the exercise when executing MODFLOW:

\begin{itemize}

\item Go to {\color{blue}Parameters} $\rightarrow$ {\color{blue} Chemistry} and use 
the {\color{blue}Write Button} to write all the input files required to run PHT3D

\item Then go to {\color{blue}Parameters} $\rightarrow$ {\color{blue}Chemistry} and use the {\color{blue}Plume Button} to start PHT3D.

\item After the end of the simulation visualise the ammonium plume and check if the 
simulation results are plausible. 

\item Plot also a breakthrough curve (BTC1) for the ammonium concentration at x = 50 m
and z = 7m. To compare the results with "observed" data go to the main menu (top row) 
to {\color{blue}Add-in}  $\rightarrow$ {\color{blue}Batch} and add the content 
of the pre-prepared python script.

\end{itemize}

Applying this little python script, which can also easily be adopted for other cases,
will show your simulated and the observed data in direct comparison:  \\

\noindent {\color{blue}
from importExport import * \\
imp =  impFile(core.gui,core) \\
f=core.fileDir+os.sep+'ex3BTC1.txt' \\
d = imp.impTabFile(f,titleNbCol=0) \\
specIn,data = d['cols'],d['data'] \\ 
tlist =core.getTlist2() \\
iper = range(len(tlist)) \\
species =['Amm','Ca','K','Mg','Na'] \\
time,conc,labs = core.onPtObs('B0',iper,'Chemistry','BTC1',species) \\
pylab.figure() \\
nrows,ncols = 3,2 \\
for i,spec in enumerate(species): \\
\indent 	pylab.subplot(nrows,ncols,i+1) \\
\indent 	col0 = specIn.index(spec) \\
\indent 	pylab.plot(data[:,0],data[:,col0],'o') \\
\indent 	pylab.plot(tlist,conc[:,i]) \\
\indent 	pylab.legend([spec]) \\
pylab.draw() \\ }


To improve the agreement between observed and simulated data we now consider ion exchange reactions. 
To include this process in the reation network and to estimate
a suitable value for the cation exchange capacity (CEC) of the exchanger sites 
proceed as follows: 

\begin{itemize}

\item To add the exchanger species to the reaction network
select go to {\color{blue}Parameters} $\rightarrow$ {\color{blue}Chemistry} 
and activate {\color{blue} X-} under the {\color{blue} Exchange Tab}.
Add a background value of 0.001 as a first estimate for the CEC. 

\item Now rerun PHT3D.

\item Once the model run is complete, reuse the python script to plot again the breakthrough 
curves (BTCs) for the ammonium, Ca and Na concentrations and repeatedly compare the results 
with the observed data.

\item Continue this manual calibration of the CEC by rerunning the model with successively improved
estimates until a good agreement is achieved.  

\end{itemize}


